% Provenance: copied from the design workspace CambridgePQ/recursion.fast, `report/sections/B-glossary.tex`, on 2026-10-19. Read-only copy; the source is authoritative.
\section{Glossary}\label{sec:glossary}

Terms are listed alphabetically. A term in \emph{italics} inside a definition has its own entry.
Section references point to the full treatment.

\begingroup\small
\begin{description}[leftmargin=0pt,itemsep=2pt,parsep=0pt,font=\normalfont\bfseries]

\item[Accepting-run bound.] The cost clause \lean{AcceptingCostAtMost}: for every index of the
  class shape, every instance and every transcript (raw prover messages and challenges), a run of
  the image that halts with \emph{accept} has \emph{cost vector} $\le c$ componentwise; the latency
  clause \lean{AcceptingLatencyAtMost} bounds \emph{work and span} the same way. Rejecting runs are
  bounded only by the \emph{fuel} (user input U2; \cref{subsec:cost-u2}).

\item[Accumulation scheme.] A scheme that defers the expensive part of verification. Each step
  checks a new instance cheaply and \emph{folds} it into a running accumulator, and a final
  \emph{decider} checks the accumulator once \cite{BCMS20,BCLMS21}. recursion.fast models it by
  three reductions, \lean{embed}, \lean{fold} and \lean{decide} (\cref{sec:recursion}).

\item[Additive FFT.] An FFT over an $\F_2$-linear subspace of a binary field, in the novel
  polynomial basis of \cite{LCH14-AdditiveFFT}. It encodes Reed--Solomon codewords on additive
  domains in $O(n\log n)$ field operations; binary fields have no large multiplicative subgroups
  of 2-power order.

\item[Admission cap.] A decidable bound in the certificate that a submission must meet in order to be
  scored: code and data at most $2^{20}$\,B, tape at most $2^{20}$\,B, $\mathit{vkLen}\le2^{10}$,
  message oracles $\ell_j\le2^{k+10}$, index oracle $\mathit{idxLen}\le2^{k+8}$, $V_k\le2^8$,
  $Q_k\le2^{16}$, at most $2^{14}$ rounds (\cref{subsec:cost-caps}).

\item[Agrees.] The verifier--spec link (\cref{def:proto-agrees}). For every class, every index of
  class shape, every instance, every raw oracle content and every challenge vector, the machine's
  decision equals the spec verifier's. Preprocessing (VK, IDX) is honest; everything the prover
  controls is universally quantified.

\item[ArkLib.] The Lean library of interactive oracle reductions and their security notions
  \cite{arklib2026}. The contract restates what it needs locally; ArkLib appears only in
  \code{Bridge/} and in the organiser library.

\item[Axiom (Lean).] A constant postulated without proof. Submissions may depend only on
  \lean{propext}, \lean{Quot.sound} and \lean{Classical.choice}; \emph{sorry} and
  \emph{native\_decide} introduce further axioms and are refused.

\item[Badge path.] The unscored harness runs that gate \emph{record} eligibility: the compiled
  indexer on hidden indices, the native prover against sampled challenges, transcript comparison
  up to $2^{16}$, measured $\le$ claimed on the trusted emulator, and prover budgets
  (\cref{subsec:arch-badge}). It runs last, as its own process; v0.3 has the hook, not the checks.

\item[BCS transformation.] The compiler from a public-coin IOP to a succinct non-interactive
  argument in the random-oracle model. It commits to oracles with Merkle trees and derives
  challenges by hashing \cite{BCS16,CY24-Book}. Its soundness follows from \emph{state-restoration
  soundness}.

\item[Capacity.] For a code of \emph{rate} $\rho$, the distance $1-\rho$, beyond which a received
  word can be close to too many codewords. \enquote{Up to capacity} proximity conjectures are
  banned in both tracks, and several have been refuted \cite{CS25-PGC,BCHKS25}.

\item[Certificate.] The \lean{Prop} structure that a submission inhabits (\lean{Certificate0},
  \lean{Certificate1}, \lean{Certificate2}; \lean{Latency1} on top of \lean{Certificate1}). Its
  fields are the shape summary, the score, security (completeness, RBR-KS, error cap), the admission
  caps, \emph{Agrees} and the cost clauses. Its scalars and the image are rendered from
  \emph{claim.json} and the image files (\cref{subsec:proto-certificate}).

\item[claim.json.] The submitted claim: per class the 18 claimed counters, the declared shape
  summary (rounds, message rounds, $E$-challenges, index challenges, total message length, tape
  input words), the image paths and the claimed scores; on the latency board also $(W_k,S_k)$
  (\cref{tab:arch-bundle}). A strict schema; every value is rendered into the theorem statement.

\item[Class.] One of the three size classes $k\in\{16,20,24\}$, with fixed shape $m=\mathrm{cols}=2^k$,
  $n_x=2^{10}$ and $\mathsf{nnz}\le2^{k+4}$. Smaller R1CS are padded to the class shape by the
  harness.

\item[Class-uniform cost theorem.] The adopted anti-overfitting rule: the cost bound is one
  theorem over \emph{every} index of the class shape, every instance and every transcript, so a
  fast path for one circuit cannot lower the score (\cref{subsec:bench-hardcode}). Not to be
  confused with the rejected \emph{flat per-instruction model}.

\item[clmul, clmulh, clmulr.] The Zbc instructions \cite{RISCV-Bitmanip}. They return,
  respectively, the low 64 bits, the high 64 bits, and bits 63--126 of the 128-bit carry-less
  product of two 64-bit registers. A $K$ multiplication is \code{clmul} plus \code{clmulh} plus a
  reduction.

\item[Comparator.] The judge for Lean proofs \cite{comparator}. It checks that each target theorem
  of the solution has exactly the type it has in the trusted challenge, and that every constant used
  by the statements is identical. It audits the axiom closure and replays the solution in a fresh
  kernel, via \emph{lean4export}. It runs first, in a fresh tree in which no solution code has run
  (Decision~A-3).

\item[CompPoly.] The Lean library of computable polynomials \cite{comppoly2026}. It provides the
  binary fields \lean{BF64} and \lean{BF64.Ext3}, a kernel-checked Rabin irreducibility
  certificate, and \lean{carryLessMul}.

\item[Completeness (perfect).] For every $(i,x,w)$ in the relation and every challenge vector, the
  honest prover's transcript is accepted (\cref{def:sec-complete}).

\item[Contract; contract id.] The trusted part of the competition: Lean definitions and statements,
  the machine, the vendored ISA files, stubs and pins. The contract id is the SHA-256 digest of the
  file that pins all protected files (\cref{subsec:arch-versioning}). \emph{Freeze} means that, after
  launch, only security fixes change the id.

\item[Cost vector.] The 18 counters \lean{CostVec}, certified per class: $n_{\mathrm{add}}$,
  $n_{\mathrm{logic}}$, $n_{\mathrm{shift}}$, $n_{\mathrm{mul}}$, $n_{\mathrm{div}}$,
  $n_{\mathrm{clmul}}$, $n_{\mathrm{ld}}$, $n_{\mathrm{ldsub}}$, $n_{\mathrm{st}}$,
  $n_{\mathrm{stsub}}$, $n_{\mathrm{br}}$, $n_{\mathrm{jmp}}$, $n_{\mathrm{sys}}$, the
  \emph{first-touch} counters $m_{\mathrm{words}}$, $m_{\mathrm{pages}}$, and the query counters
  $q^{\mathrm{count}}$, $q^{\mathrm{lev}}$, $q^{\mathrm{bytes}}$. Each \emph{zkVM profile} prices it;
  re-weighting needs no re-proof (\cref{subsec:cost-vector}).

\item[Cycle.] One retired instruction of the contract machine (a step, against the \emph{fuel}).
  On the recursion board its price is set by its class in the \emph{cost vector}; on the latency
  board the pipeline timer assigns it cycles.

\item[decide +kernel.] Proving a decidable proposition by kernel reduction alone. It trusts no
  compiled code, but it can be very slow and memory-hungry; compare \emph{native\_decide}.

\item[Decider.] The final check of an \emph{accumulation scheme}, \lean{decide}: $\Racc\to\top$. It
  enters $\mathrm{Score}_{T2}$ as $S_{\mathrm{dec}}/64$ and is bounded only by the fuel; an all-runs
  bound is planned, not in v0.3.

\item[Definition hole.] A definition whose body is supplied by the submission, such as the
  submission bundle, and which comparator matches by name and type only. The certificate must pin it
  so tightly that any inhabitant is a legitimate entry.

\item[ECALL.] The RISC-V environment-call instruction. In the contract, register \code{t0} selects
  the system call: \code{0 HALT}, \code{1 QUERY}, \code{2 HASH} (reserved for the ROM track; traps),
  \code{3 SPAWN}, \code{4 JOIN} (\emph{harts} only).

\item[Elaboration vs replay.] Elaboration turns source into kernel terms by running tactics and
  metaprograms; it parallelises across modules. Replay re-checks the exported declarations in a
  fresh kernel; it is single-threaded and sets a floor on check time (\cref{subsec:arch-budgets}).

\item[Environment gate.] The harness check, run on the elaborated \code{.olean}s in a second fresh
  tree after comparator, that rejects \lean{sorryAx} and non-whitelisted axioms in the certificate's
  closure, \lean{unsafe}/\lean{partial}/\lean{opaque} untrusted constants, any
  \lean{implemented\_by}, \lean{extern}, \lean{csimp}, \lean{init}, \lean{builtin\_init} or
  \lean{export} entry in an untrusted module, and badge-path fields (\lean{VK}, \lean{IDX}, prover,
  extractor) without compiled code. A module
  is trusted iff its files match the \emph{organiser manifest}. \lean{noncomputable} means
  \enquote{no compiled code}, not \enquote{infeasible}, so the gate does not establish feasibility
  (\cref{subsec:arch-ci}).

\item[Error cap; round cap.] Every per-round RBR-KS error is at most $2^{-128}$, and every
  protocol or reduction has at most $\mathit{roundCap}=2^{14}$ rounds (\cref{subsec:sec-cap}).

\item[Extractor.] The algorithm that recovers a witness from an accepting prover. In the contract
  it is \emph{straight-line} (it reads the oracles and never rewinds), named and computable.
  Without a running-time bound it may brute-force the witness, so L1 \enquote{knowledge} equals
  RBR soundness (\cref{subsec:sec-extractor}).

\item[Fiat--Shamir.] Replacing the verifier's random challenges by hashes of the transcript so far
  \cite{FS86}; its multi-round security is analysed in \cite{CCHLRRW19,AFK22}. recursion.fast fixes the transcript prefix (statement binding) and
  charges a static surcharge $\mathit{FS}_k$ for the in-circuit hashing (\cref{subsec:sec-fs}).

\item[First touch.] The first access to a doubleword or a 1\,KiB page in a run; the counters
  $m_{\mathrm{words}}$, $m_{\mathrm{pages}}$ count them (code and data apart, a written page twice),
  pricing memory in profiles that charge paging.

\item[Flat per-instruction model.] One fixed weight per instruction class for every prover, as in
  sig.golf and recursion.fast v0.2 (kept as a display column). Rejected because agents overfit it
  (user input U3; \cref{subsec:cost-diagnosis}).

\item[Folding scheme.] An accumulation scheme in which two instances of a relation are reduced to
  one instance of the same relation, as in Nova \cite{KST22-Nova}.

\item[Fuel.] The step limit $F=2^{36}$ per hart of \lean{execute}. Running out of fuel counts as
  rejection, so no separate termination theorem is needed. A child hart runs on its parent's
  remaining fuel, so a run with harts may take $(\#\mathrm{spawns}+1)\cdot2^{36}$ steps.

\item[fullRecAcc.] The T2 baseline: \lean{embed} is the whole verifier, and \lean{fold} and
  \lean{decide} accept everything over $\Racc=\mathrm{univ}$. It scores exactly the T1 cost.

\item[Grinding.] A proof of work by the prover before a challenge is derived, which raises the cost
  of Fiat--Shamir attacks and lets the protocol use fewer queries. It exists only after
  compilation, so it is outside the IOP model (ROM track).

\item[Hart; SPAWN, JOIN.] A hardware thread. \code{SPAWN} runs a child hart with declared read and
  write regions, sequentialised in the model; \code{JOIN} makes its writes visible to the parent. Only
  images with the \lean{harts} flag may use them; races reject (\cref{subsec:proto-machine}).

\item[Harvard memory.] Code and data in separate address spaces. Code cannot be read or written as
  data. The data memory is flat, $2^{36}$ bytes, naturally aligned.

\item[Holography; indexer.] An \emph{indexer} preprocesses the index $i$ (the R1CS matrices) once
  into a short verification key $\vk\,i$ and an index oracle $\idx\,i$. The verifier reads the key
  and queries the oracle, so it never reads $i$ \cite{CHMMVW20-Marlin,COS20}
  (\cref{def:bg-holographic}).

\item[IOP; IOPP.] An interactive oracle proof \cite{BCS16,RRR16}: the prover sends oracles, and the
  verifier sends random challenges and reads a few positions of the oracles. An IOPP is an IOP of
  proximity, which tests that an oracle is close to a code. FRI, STIR and WHIR are IOPPs.

\item[IVC; PCD.] Incrementally verifiable computation \cite{Val08} and its generalisation to DAGs,
  proof-carrying data \cite{CT10}: each step proves that the previous proofs verified. They are
  built from recursive SNARKs or accumulation.

\item[Johnson regime.] Proximity parameters $\delta<1-\sqrt\rho$ (the Johnson bound), within which
  list decoding of RS codes has polynomially small lists. Proven proximity gaps cover it
  \cite{BCIKS20,BCHKS25}. ArkLib formalises the result for lines only, so in practice WHIR-type
  designs in this regime need the named hypothesis \lean{McaJohnson} (\cref{subsec:sec-hypotheses}).

\item[$K$, $E$.] The relation field $K=\mathrm{GF}(2^{64})=\F_2[x]/(x^{64}+x^4+x^3+x+1)$, with
  8-byte elements, and the challenge field $E=K[y]/(y^3+y+1)=\mathrm{GF}(2^{192})$, with 24-byte
  elements (\cref{subsec:proto-fields}).

\item[Kernel (Lean).] The small type checker that accepts or rejects declarations. It is the root of
  trust. nanoda \cite{nanoda} is an independent implementation and the second kernel of the
  harness; \code{lean4checker} \cite{lean4checker} replays in Lean's own kernel, a further replay
  but not kernel diversity.

\item[KSF.] Knowledge state function: predicates $S_m(s,\mathrm{tr},w)$, together with the
  per-move backward extractor, satisfying three laws (empty, next, full). A challenge is bad if it
  makes some $w$ valid when its pullback was not (\cref{subsec:sec-rbrks}).

\item[Latency board; Latency1.] The secondary leaderboard for native multicore latency: a proved
  pipeline \emph{work} $W_k$ and \emph{span} $S_k$ per class, scored by
  $\mathrm{Score}_{\mathrm{lat}}$ (\cref{subsec:cost-lat}). Its certificate clause is
  \lean{Latency1}, rendered from the stub \code{L1.lean.in}; the file name does not refer to
  security level L1.

\item[Level (L0--L4).] A security level of the contract (\cref{sec:protocol}). L0 is the baseline;
  L1 is the v1 base track with one-step accumulation; L2 adds the tree theorem, efficient
  extraction and named hypotheses; L3 is the ROM argument; L4 is full recursion, as a target
  schema. Not to be confused with the library components lib-L0--lib-L6 (\cref{sec:issues}) or
  with the latency stub \code{L1.lean.in}.

\item[lean4export.] The tool that exports Lean declarations in a plain-text format, for replay by an
  independent checker \cite{lean4export}.

\item[leanerVM; leanVM.] leanVM is a minimal zkVM over binary fields, with 2-to-1 recursion
  \cite{leanVM}. leanerVM is its Lean formalisation, the source of the reused components,
  \lean{McaJohnson} and the toolchain pins.

\item[List decoding.] Decoding beyond half the minimum distance: the output is the list of all
  codewords within distance $\delta$. Soundness analyses in this regime pay for the list size, and
  extractors need efficient list decoders.

\item[Mathlib.] Lean's community mathematics library. Besides Lean core and the vendored RISC-V
  model, it is the only import of the trusted contract.

\item[MCA.] Mutual correlated agreement: if a random linear combination of words is close to the
  code, then all the words agree with codewords on a \emph{common} large set. It is the
  proximity-gap form used by WHIR-type protocols \cite{ACFY24-WHIR,BCGM25-MCA}.

\item[native\_decide.] Proving a decidable proposition by running compiled code. It trusts the compiler and
  adds the axiom \lean{Lean.ofReduceBool}, so it is refused.

\item[Negative and positive controls.] Kernel-checked tests of the contract. A negative control
  (for instance an accept-all verifier) must fail the certificate; a positive control (for instance
  the baseline) must satisfy it (\cref{tab:sec-controls}).

\item[OOD sample.] An out-of-domain query: the verifier asks for the evaluation of the committed
  polynomial at a random point outside the evaluation domain. In the list-decoding regime this pins
  down a unique codeword (DEEP, STIR, WHIR).

\item[Oracle; leaf.] Oracle $j=0$ is $\idx\,i$, and oracles $j\ge1$ are the prover's messages. A
  \code{QUERY} returns one leaf of $\mathit{leaf}_j$ elements, priced as an opening of a Merkle path
  of depth $\log_2(\ell_j/\mathit{leaf}_j)$.

\item[Organiser package; manifest.] The organiser's build of the contract, vendored model, library,
  controls, bridge, \code{gate} and \code{axiomaudit}, made outside every tree in which solution
  code runs and mounted read-only. A committed SHA-256 manifest of every toolchain, package and
  organiser \code{.olean} is checked before and after each harness stage (\cref{subsec:arch-ci}).

\item[pad.] The contract maps \lean{SpMat.pad}, \lean{PubInput.pad} that embed a smaller R1CS into
  the class shape (entries moved by a column map, zero public inputs). The library theorem
  \lean{pad\_sat\_iff} proves that padding preserves and reflects satisfiability.

\item[PCS.] Polynomial commitment scheme. Hash-based PCSs combine an RS code, Merkle commitments and
  an IOPP. A knowledge-sound PCS over binary RS codes is the critical path of the organiser library.

\item[PIOP.] Polynomial IOP: an IOP whose oracles are low-degree polynomials, compiled to an IOP over
  codes by a PCS.

\item[Precompile.] An instruction or system call with a fixed price that performs a heavy function
  natively, such as a hash or a field multiplication. v1 has none beyond Zbc; \code{HASH} and
  \code{EMUL} are reserved for later boards.

\item[Proximity gap.] For an affine family of words and a code, either almost every member is
  $\delta$-close to the code, or almost none is. The exceptional fraction is the error that enters
  the soundness of FRI- and WHIR-type tests \cite{BCIKS20} (\cref{subsec:bg-proximity}).

\item[QProg.] The contract's query program, a free monad over \code{QUERY} with pure local
  computation. It is the verifier form of T0 and the midpoint of the T1 refinement
  (\cref{subsec:proto-qprog}).

\item[QUERY.] System call 1: \code{a0} is the oracle $j$, \code{a1} the leaf index and \code{a2} the
  destination. It writes the leaf to memory, and an out-of-range request traps.

\item[R1CS.] Rank-1 constraint system: $(Az)\circ(Bz)=Cz$ with $z=(1,x,w)$ over $K$
  (\cref{def:bg-r1cs}).

\item[Rate.] $\rho=\dim/\text{length}$ of a code. A lower rate gives more distance per query but
  longer oracles, so a higher prover cost.

\item[RBR soundness.] Round-by-round soundness \cite{CCHLRRW19}: there is a state function under
  which a false statement starts doomed, the verifier rejects every full transcript whose state is
  doomed, and from every doomed prefix the next challenge leaves the doomed set with probability
  at most $\eps_i$
  (\cref{def:bg-rbr}).

\item[RBR-KS.] Round-by-round knowledge soundness: the knowledge analogue of RBR soundness, stated
  with a \emph{KSF} and a per-move backward extractor. The contract uses the worst-case counting
  form $\#\mathrm{bad}_i\le\eps_i|C_i|$ \cite{BCFW25-WARP} (\cref{subsec:sec-rbrks}).

\item[Record.] A leaderboard entry under a given contract id: certified, and passed the badge
  path. Records are revalidated on every toolchain or semantics change.

\item[Recursion board.] The primary leaderboard: the certified \emph{cost vector} priced under four
  \emph{zkVM profiles}, scored by $\mathrm{Score}_{\mathrm{zk}}$ (\cref{subsec:cost-zk}).

\item[RecursionHeuristic.] The named deployment assumption that instantiating the random oracle
  inside the recursion preserves straight-line extraction. It is never a hypothesis of a scored
  theorem (\cref{subsec:rec-cannot}).

\item[Ring switching.] A technique that reduces an evaluation claim about a multilinear polynomial
  over a small binary field, with packed coefficients, to a claim over an extension field. It
  allows small-field commitments with large-field challenges \cite{DP24-FRIBinius}.

\item[RS code.] Reed--Solomon code: evaluations of polynomials of degree $<\rho n$ on a domain of
  size $n$. In recursion.fast the domains are additive subspaces of $K$.

\item[riscv-zkvm.] The Lean RISC-V model whose \code{Rv64/*} files are vendored byte-identically
  from commit \repo{6810fc1} \cite{riscvzkvm2026}.

\item[ROM track.] A planned track in the random-oracle model, with a \code{HASH} system call, real
  Merkle and Fiat--Shamir costs, and grinding (security level L3).

\item[Round shape.] A contract-fixed round type: \lean{msg o}, a raw vector
  \lean{Fin (2\^{}o.logLen) → K} from the prover, or \lean{chal nE idx}, $n_E$ elements of $E$
  plus index challenges from the verifier (\cref{subsec:proto-rounds}).

\item[RV64IM; W-ops.] The 64-bit RISC-V base integer ISA with the multiply/divide extension. The
  W-ops (\code{ADDW}, \code{MULW}, \code{SRAW} and so on) operate on 32-bit words and sign-extend
  the result.

\item[Score.] Per profile $P$ and class $k$,
  \[S_{P,k}=\bigl(w_P\cdot c_k+(P_{H,P}+35)q^{\mathrm{lev}}_k+P_{H,P}(q^{\mathrm{bytes}}_k/64+q^{\mathrm{count}}_k)
  +\mathit{FS}_{P,k}+\lceil\lambda_PB_k\rceil\bigr)\max(1,V_k/16)^{1/4};\]
  $\mathrm{Score}_{\mathrm{zk}}=(\prod_P\prod_kS_{P,k})^{1/12}$;
  $\mathrm{Score}_{T2}$ uses $S_{\mathrm{leaf}}+S_{\mathrm{fold}}+S_{\mathrm{dec}}/64$;
  $\mathrm{Score}_{\mathrm{lat}}=(\prod_k(W_k/16+S_k))^{1/3}$. Lower is better. The certificate
  proves the claimed fourth powers (\lean{scoreZk4}) and the scorer takes the root
  (\cref{subsec:cost-score}).

\item[sorry.] Lean's placeholder for a missing proof. It adds the axiom \lean{sorryAx} and is
  refused.

\item[State-restoration soundness.] Soundness against provers that may rewind the verifier to any
  earlier state and draw fresh challenges, up to $t$ times. It governs the security of BCS and
  Fiat--Shamir in the random-oracle model, and RBR soundness implies it \cite{CY24-Book,Hol19}.

\item[Submission root.] The single flat directory per pull request, holding the image,
  \emph{claim.json} and the Lean solution. It is exported at blob level from a frozen commit.

\item[Sumcheck.] The interactive protocol that reduces a claim about the sum of a multivariate
  polynomial over the Boolean hypercube to one evaluation, with error $d/|E|$ per round.

\item[Track (T0--T3).] A score track. T0 is Lean-only, with a \emph{QProg} verifier and its three
  query counters. T1 is RISC-V images on the \emph{recursion board} (and optionally the
  \emph{latency board}). T2 is recursion: leaf, fold and decider. T3 is arithmetised recursion,
  future work.

\item[Tree composition theorem.] \lean{tree\_secure}: an accumulation tree of $N$ leaves is complete
  and RBR-KS, with $k_{\mathrm{tree}}$ rounds and independent per-node challenges. It is an
  organiser obligation at L2.

\item[Trusted surface.] Everything a reader must accept for a record to mean what it says: the
  contract text, the vendored ISA, Lean's kernel, elaborator, Lake and compiler, comparator and
  nanoda, the harness and its host tools, and the cited meta-theorems (\cref{tab:arch-trusted}).

\item[Unique decoding.] Decoding within half the minimum distance, $\delta<(1-\rho)/2$ for RS
  codes, where the closest codeword is unique. It is the regime of the unconditional track: its
  proximity gaps are proved in ArkLib \cite{BCIKS20}.

\item[$V_k$ (blow-up).] $\sum_{j\ge1}\ell_j/(N_k+2^{k+4})$, the declared prover-message length
  relative to the instance size. It enters the score through the factor $\max(1,V_k/16)^{1/4}$ and
  is capped at $2^8$.

\item[VCVio.] The Lean framework for oracle computations and cryptographic proofs
  \cite{tuma2026vcvio}, on which ArkLib builds.

\item[W, S (work, span).] Latency-board quantities: $W$ is the pipeline time of all harts plus a
  fixed charge per spawn, $S$ the root hart's clock, which waits at each \code{JOIN}.

\item[Zbc.] The RISC-V carry-less multiplication extension (\emph{clmul, clmulh, clmulr}), with its
  own counter, priced per profile (0.97--1.38; 4 in the v0.2 display column).

\item[zkVM profile.] A weight vector on the counters plus a hash price $P_{H,P}$ and a byte charge,
  calibrated from a zkVM's published cost tables: ZisK, SP1-Hypercube, RISC~Zero, Jolt
  (\cref{tab:cost-profiles}).

\end{description}
\endgroup
